% October 9, 2026 build (agent draft): conclusion following the deck's
% "Conclusion" frame, with the property-tax claim stated as the Oct 9
% rebate-held run supports it (the deck's bullet says the tax moves space
% toward young families; rent per room does not fall in any rebated arm).
\section{Conclusion}
\label{sec:conclusion}

I build a life-cycle model of fertility with realistic housing, in which
households choose when to have children, how much space to occupy, and
whether to rent or own it. Children need space, the large homes must be
bought with a down payment, and the old hold most of them. I calibrate the
model to the United States in 2007 and treat the fall in fertility since then
as the start of a transition to a smaller population, along which a durable
housing stock makes space cheaper for the smaller cohorts that follow.

Three results emerge. First, the per-period cost of the space a child
requires, and the buffer a household holds against it, set the number of
children; the down payment, credit and the house price at a given rent set
their timing. Second, mortgage credit raises the ownership of the young by
20 to 25 percentage points and leaves completed fertility within about one
percent, except where family-sized homes cannot be rented. Third, a property
tax rebated equally to all households raises births when a durable stock lets
the price absorb the tax, by 3 to 4 percent, and lowers them when supply
passes the tax into rents. Its effect runs through the rebate: it moves income,
not space, toward young families, paid for by the owners at the time of the
reform.

The analysis has limitations that future work should address. The
preference shocks that account for the decline in fertility since 2007 fall
on every cohort alive, so the model overstates childlessness and misses the
postponement of first births, and with only preference shocks it predicts a
falling house price where the data show a rising one. The model also gives
young households too many homes and the old too few large ones. A housing
shock that matches the observed rise in prices and rents, and a fertility
shock that falls on young cohorts, are the natural next steps.
